%%
%% Arquivo principal adaptado para:
%% - Trabalho de Conclusão de Curso de Graduação - Eng. Computação e Mecatrônica - UFRN
%%
%% NOTA: A PUBLICAÇÃO DESTE MODELO VISA APENAS ORIENTAR OS PÓS-GRADUANDOS
%% NA PREPARAÇÃO DE SEUS TEXTOS. O PPgEE DA UFRN NÃO PROVÊ ASSISTÊNCIA NO
%% USO DAS FERRAMENTAS NECESSÁRIAS AO USO DESTE MODELO (LATEX, XFIG, ETC.)
%%
%% Baseado no modelo prévio para Tese de Pós Graduação PPGEEC do Prof. Dr.Adelardo Medeiros, dezembro de 2005.
%% Revisado pelos alunos de Metodologia da Pesquisa Científica de 2016.1.


%% -------------------------------------------------------------------------
%%
%% NOTA: Modelo adaptado em Julho de 2022 para template de TCC em Eng. de Computação 
%% e Mecatrônica da UFRN/CT.
%%
%% Válber César Cavalcanti Roza (Sec. de Coordenação de Eng.Mecatrônica CTEC/UFRN),20 de Julho, 2022.
%% Revisado pelo secretariado das Coordenações dos respectivos cursos.
%%
%% -------------------------------------------------------------------------


% DEFINIÇÕES GLOBAIS

% Esta primeira linha dá informações gerais sobre o documento.
% PARA A VERSÃO FINAL:
% papel A4, letra grande (12pt), openright (capítulos só iniciam em
% página direita, se necessário incluindo uma página em branco),
% twoside (o documento vai ser impresso em frente e costa)
\documentclass[a4paper,12pt,openright,twoside]{book}
% PARA A QUALIFICAÇÃO E PARA A VERSÃO INICIAL:
% papel A4, letra grande (12pt), openany (capítulos iniciam em
% qualquer página), oneside (o documento vai ser impresso só na frente)
%\documentclass[a4paper,12pt,openany,oneside]{book}


% Use estes pacotes para poder digitar diretamente as letras com acento
% e para que a hifenização funcione corretamente
\usepackage[utf8]{inputenc}
\usepackage{ae}
% Para usar fontes standard ao invés das do LaTeX (gera melhores PDFs)
\usepackage{pslatex}
% Para a hifenização em português
\usepackage[portuges, brazil]{babel}
% Para que os primeiros parágrafos das seções também sejam indentados
\usepackage{indentfirst}
% Para poder incluir gráficos (figuras)
\usepackage{graphicx}
% Para poder fazer glossário ou lista de símbolos
% Use a segunda opção se quiser incluir na definição do símbolo a
% página e/ou a equação onde ela foi definida
\usepackage[portuguese,noprefix]{nomencl}
%\usepackage[portuguese,noprefix,refeq,refpage]{nomencl}
% Para permitir espaçamento simples, 1 1/2 e duplo
\usepackage{setspace}
% Para usar alguns comandos matemáticos avançados muito úteis
\usepackage{amsmath}
% Para poder usar o ambiente "comment"
\usepackage{verbatim}
% Para poder ter tabelas com colunas de largura auto-ajustável
\usepackage{tabularx}
% Para executar um comando depois do fim da página corrente
\usepackage{afterpage}
% Para formatar URLs (endereços da Web)
\usepackage{url}
% Para reduzir os espaços entre os ítens (itemize, enumerate, etc.)
% Este pacote não faz parte da distribuição padrão do LaTeX.
\usepackage{lib/noitemsep}
% Para as citações bibliográficas
%\usepackage[alf]{abntex2cite}
\usepackage[abbr]{lib/harvard}	% As chamadas são sempre abreviadas
%\harvardparenthesis{square}	% Colchetes nas chamadas
%\harvardyearparenthesis{round}	% Parêntesis nos anos das referências
%\renewcommand{\harvardand}{e}	% Substituir "&" por "e" nas referências

% PACOTES OPCIONAIS

% Para poder incluir arquivos Postscript com cores (do Xfig, por exemplo)
\usepackage{color}
% Para ter células em tabelas que ocupam mais de uma linha
\usepackage{multirow}
% Para poder ter tabelas longas em mais de uma página
%\usepackage{longtable}
% Para poder escrever partes do texto em "n" colunas
%\usepackage{multicol}
% Se você quiser personalizar os cabeçalhos das páginas
%\usepackage{fancyheadings}
% Para incluir algoritmos e listagens de códigos
%\usepackage{listings}
% Para incluir pseudocódigos
\usepackage[portuguese, ruled, linesnumbered]{algorithm2e}
% Capítulos com títulos em um formato "decorado"
\usepackage{lib/capitulos}
% Hiperreferências dentro do texto e montagem dos links do índices dos
% para os leitores de pdf (deve ser o último pacote a ser inserido).
\usepackage[breaklinks]{hyperref}
% Referência correta de ambientes flutuantes (como figuras, tabelas e algoritmos).
\usepackage[all]{hypcap}

% NOVOS COMANDOS

% As definições dos novos comandos estão agrupadas no arquivo "comandos.tex"
% Esta separação é opcional: se você preferir, pode por as definições
% diretamente neste arquivo
% newcommand define novos comandos, que podem passar a ser usados da
% mesma forma que os comandos LaTeX de base.

% Implicação em fórmulas
\newcommand{\implica}{\quad\Rightarrow\quad} %Meio de linha
\newcommand{\implicafim}{\quad\Rightarrow}   %Fim de linha
\newcommand{\tende}{\rightarrow}
\newcommand{\BibTeX}{\textsc{B\hspace{-0.1em}i\hspace{-0.1em}b\hspace{-0.3em}}\TeX}

% Fração com parentesis
\newcommand{\pfrac}[2]{\left(\frac{#1}{#2}\right)}

% Transformada de Laplace e transformada Z
%\newcommand{\lapl}{\makebox[0cm][l]{\hspace{0.1em}\raisebox{0.25ex}{-}}\mathcal{L}}
\newcommand{\lapl}{\pounds}
\newcommand{\transfz}{\mathcal{Z}}

% Não aparecer o número na primeira página dos capítulos
\newcommand{\mychapter}[1]{\chapter{#1}\thispagestyle{empty}}

% Os capítulos sem número
\newcommand{\mychapterast}[1]{\chapter*{#1}\thispagestyle{empty}
\chaptermark{#1}
\afterpage{\markboth{\uppercase{#1}}{\rightmark}}
\markboth{\uppercase{#1}}{}
}

% Seções sem número
\newcommand{\mysectionast}[1]{\section*{#1}
\addcontentsline{toc}{section}{#1}
\markright{\uppercase{#1}}
}

% No tabularx, as celulas devem ser centradas verticalmente
\renewcommand{\tabularxcolumn}[1]{m{#1}}

% Células centralizadas horizontalmente no tabularx
\newcolumntype{C}{>{\centering\arraybackslash}X}

%% Abrevia figuras e tabelas
%\def\figurename{Fig.}
%\def\tablename{Tab.}


%
% As margens
%

% Direção horizontal

% Margem interna
% Nas páginas ímpares
\setlength{\oddsidemargin}{3.5cm}         % Margem real desejada
% Nas páginas pares
\setlength{\evensidemargin}{2.5cm}        % Margem real desejada
% Largura do texto
\setlength{\textwidth}{15cm}
% As margens laterais no LaTeX são sempre 1 polegada maiores do que as
% fixadas. Se foi fixada \setlength{\oddsidemargin}{3.5cm}, a margem
% real seria de 3.5+2.54=6.04cm. Para permitir que você não tenha que
% fazer esta conta, pode usar o número desejado nas linhas anteriores
% e a gente subtrai 1in nas próximas linhas
\addtolength{\oddsidemargin}{-1in}
\addtolength{\evensidemargin}{-1in}
% Note que a margem direita não é fixada diretamente:
% ela é obtida subtraindo-se os outros valores da largura da página.
% 3.5+15+x=21cm (largura A4) -> x = margem externa = 2.5cm

% Direção vertical

% Margem superior (entre o topo da folha e o cabeçalho), altura do
% cabeçalho e distância entre o fim do cabeçalho e o início do texto
\setlength{\topmargin}{2.0cm}             % Margem real desejada
\setlength{\headheight}{1.0cm}
\setlength{\headsep}{1.0cm}
% Altura do texto (sem cabeçalho e rodapé)
\setlength{\textheight}{22.7cm}
% Distância do fim do texto ao rodapé
\setlength{\footskip}{1.0cm}
% A margem superior no LaTeX é sempre 1 polegada maior do que a
% fixada. Se foi fixada \setlength{\topmargin}{2.0cm}, a margem
%real seria de 2.0+2.54=4.54cm. Para permitir que você não tenha que
% fazer esta conta, pode usar o número desejado na linha anterior
% e a gente subtrai 1in na próxima linha
\addtolength{\topmargin}{-1in}
% Note que a margem inferior não é fixada diretamente:
% ela é obtida subtraindo-se os outros valores, sem incluir o
% "footskip", da altura da página.
% 2.0+1.0+1.0+22.7+x=29.7cm (altura A4) -> x = margem inferior = 3cm

%
% O estilo das referências bibliográficas
%

\bibliographystyle{bibliografia/ppgee}

%
% O espaçamento entre linhas
%

% As páginas iniciais são sempre em espaçamento simples
\singlespacing

% Para a criação do glossário (ou lista de símbolos)
\makenomenclature

% Lista de arquivos a serem processados. Estes comandos "includeonly" são
% dispensáveis e devem obrigatoriamente ser comentados na hora de gerar o
% documento definitivo. Eles servem para compilar apenas parte do documento.
% É útil, durante a redação, para que não se tenha de compilar todo o
% documento a cada vez que se faz uma alteração. Por exemplo, se eu estou
% fazendo alterações na dedicatória e as outras partes não têm interesse no
% momento, posso incluir (descomentar) a linha "\includeonly{preambulo}"
%\includeonly{rosto}
%\includeonly{catalograficos}
%\includeonly{preambulo}
%\includeonly{resumos}
%\includeonly{introducao/introducao}
%\includeonly{estilo/estilo}
%\includeonly{matematica/matematica}
%\includeonly{figuras/figuras}
%\includeonly{conclusao/conclusao}
%\includeonly{apendice/apendice}

% Inicia o texto
\begin{document}

% Paginas iniciais (sem numeração)
\pagestyle{empty}

% Página de rosto (capa interna)
%
% ********** Página de Rosto
%

% titlepage gera páginas sem numeração
\begin{titlepage}

\begin{center}

\small

\begin{tabularx}{\linewidth}{@{}l@{}C@{}r@{}}
\parbox[c]{3cm}{\includegraphics[width=\linewidth]{./figuras/UFRN}} &
\begin{center}
\textsf{\textsc{Universidade Federal do Rio Grande do Norte\\
Centro de Tecnologia\\
Graduação em Engenharia de Computação}}
\end{center} &
\parbox[c]{3cm}{\includegraphics[width=\linewidth]{./figuras/dca_logo.png}}
\end{tabularx}

\vfill
\LARGE
\textbf{Nurse}
\vfill
\Large
\textbf{Dorgival da Rocha Filho}
\vfill

\normalsize

Orientador: Prof. Dr. Severino Lampeão
\\[2ex] Co-orientador: Prof. Dr. Zé Baiano % Opcional

\vfill

\hfill


\vfill

\large

Natal, RN, \today

\end{center}

\end{titlepage}

% Página de rosto (capa interna)
%
% ********** Página de Rosto
%

% titlepage gera páginas sem numeração
\begin{titlepage}

\begin{center}

\small

\begin{tabularx}{\linewidth}{@{}l@{}C@{}r@{}}
\parbox[c]{3cm}{\includegraphics[width=\linewidth]{./figuras/UFRN}} &
\begin{center}
\textsf{\textsc{Universidade Federal do Rio Grande do Norte\\
Centro de Tecnologia\\
Graduação em Engenharia de Computação}}
\end{center} &
\parbox[c]{3cm}{\includegraphics[width=\linewidth]{./figuras/dca_logo.png}}
\end{tabularx}

\vfill
\LARGE
\textbf{Nurse: aplicação mobile para controle dos dados de vacinação e aumento de produtividade de profissionais da saúde}
\vfill
\Large
\textbf{Dorgival da Rocha Filho}
\vfill

\normalsize

Orientador: Prof. Dr. Itamir de Morais Barroca Filho

\vfill

\hfill
\parbox{0.5\linewidth}{\textbf{%
Trabalho de Conclusão de Curso de Graduação} na modalidade Monografia, 
submetido como parte dos requisitos necessários para 
conclusão do curso de Engenharia de Computação pela 
Universidade Federal do Rio Grande do Norte (UFRN/CT).
}

\vfill

\large

Natal, RN, \today

\end{center}

\end{titlepage}


% Ficha catalográfica: os dados catalográficos devem ser fornecidos
% pela BCZM.
% Só são incluídos na versão final da tese ou dissertação. Não são
% incluídos nem na proposta de tema de qualificação nem na versão
% preliminar da tese ou dissertação: nestes casos, comente a próxima linha.
%
% ********** Ficha Catalográfica
%

\newpage

\begin{center}
  % Aqui não se usou \vfill porque o \vfill é construído internamente com
  % o comando \vspace. Espaços verticais no início da folha com \vspace
  % são ignorados. Para que isto não ocorra deve-se usar o \vspace*
  % \vspace*{\fill} é como se fosse um \vfill*
  \vspace*{\fill}
  \fontsize{10}{12}\selectfont
  
  \textsf{Universidade Federal do Rio Grande do Norte - UFRN \\
  Sistema de Bibliotecas - SISBI \\
  Catalogação de Publicação na Fonte. UFRN - Biblioteca Central Zila Mamede}
  
  \vspace{2ex}
  
  \begin{minipage}[c]{12.5cm}
    \begin{tabular}{|p{\linewidth}|} \hline \texttt{Rocha Filho, Dorgival da.} \linebreak
    \texttt{\hspace*{1em} Nurse: aplicação mobile para controle dos dados de vacinação \linebreak
    e aumento de produtividade de profissionais da saúde / Dorgival \linebreak
    da Rocha Filho -- Natal, 2022.} \linebreak
    \hspace*{1em} \texttt{88 f.: il}
    \\[1ex]
    \texttt{\hspace{1em} Monografia (graduação) - Universidade Federal do Rio Grande \linebreak
    do Norte, Centro de Tecnologia, Curso de Engenharia de \linebreak Computação.} \linebreak
    \hspace*{1em} \texttt{Orientador: Prof. Dr. Itamir de Morais Barroca Filho.}
    \\[3ex]
    \texttt{\hspace{1em} 1. Framework Flutter - Monografia. 2. Vacinação - Monografia. \linebreak 3. Aplicação Móvel - Monografia.
    I. Barroca Filho, Itamir de \linebreak Morais.
    II. Título.}
    \\
    \texttt{RN/UF/BCZM} \hfill \texttt{CDU 004.4:614.47} \hspace{3cm} \\
    \\
    \hline
    \end{tabular} 
    \\[3ex]
    \texttt{Elaborado por Fernanda de Medeiros Ferreira Aquino - CRB-15/301}
  \normalsize
  \end{minipage}
\end{center}
\vspace{4cm}


% Assinaturas da banca, dedicatória e agradecimentos
% Só são incluídos na versão final da tese ou dissertação. Não são
% incluídos nem na proposta de tema de qualificação nem na versão
% preliminar da tese ou dissertação: nestes casos, comente a próxima linha.
%
% ********** Página de assinaturas
%

\begin{titlepage}

\begin{center}

\LARGE

\textbf{Nurse: aplicação mobile para controle dos dados de vacinação e aumento de produtividade de profissionais da saúde}

\vfill

\Large

\textbf{Dorgival da Rocha Filho}

\end{center}

\vfill

% O \noindent é para eliminar a tabulação inicial que normalmente é
% colocada na primeira frase dos parágrafos
\noindent
% Descomente a opção que se aplica ao seu caso
% Note que propostas de tema de qualificação nunca têm preâmbulo.
Monografia aprovada em 09 de dezembro de 2022, pela banca examinadora composta
pelos seguintes membros:

% Os nomes dos membros da banca examinadora devem ser listados
% na seguinte ordem: orientador, co-orientador (caso haja),
% examinadores externos, examinadores internos. Dentro de uma mesma
% categoria, por ordem alfabética

\begin{center}

\vspace{1.5cm}\rule{0.95\linewidth}{1pt}
\parbox{0.9\linewidth}{%
Prof. Dr. Itamir de Morais Barroca Filho (orientador) \dotfill\ DTI-IMD/UFRN}

% \vspace{1.5cm}\rule{0.95\linewidth}{1pt}
% \parbox{0.9\linewidth}{%
% Prof. Dr. YYYYY (co-orientador) \dotfill\ MCA/UFRN}

\vspace{1.5cm}\rule{0.95\linewidth}{1pt}
\parbox{0.9\linewidth}{%
Prof. Dr. André Morais Gurgel \dotfill\ DEPAD/CCSA/UFRN}

\vspace{1.5cm}\rule{0.95\linewidth}{1pt}
\parbox{0.9\linewidth}{%
Prof. Dr. Jean Mário Moreira de Lima \dotfill\ IMD/UFRN}

\end{center}

\end{titlepage}

%
% ********** Dedicatória
%

% A dedicatória não é obrigatória. Se você tem alguém ou algo que teve
% uma importância fundamental ao longo do seu curso, pode dedicar a ele(a)
% este trabalho. Geralmente não se faz dedicatória a várias pessoas: para
% isso existe a seção de agradecimentos.
% Se não quiser dedicatória, basta excluir o texto entre
% \begin{titlepage} e \end{titlepage}

\begin{titlepage}

\vspace*{\fill}

\hfill
\begin{minipage}{0.5\linewidth}
\begin{flushright}
\large\it
Aos meus ......
\end{flushright}
\end{minipage}

\vspace*{\fill}

\end{titlepage}

%
% ********** Agradecimentos
%

% Os agradecimentos não são obrigatórios. Se existem pessoas que lhe
% ajudaram ao longo do seu curso, pode incluir um agradecimento.
% Se não quiser agradecimentos, basta excluir o texto após \chapter*{...}

\chapter*{Agradecimentos}
\thispagestyle{empty}

\begin{trivlist}  \itemsep 2ex

\item Ao meu orientador e ao meu co-orientador, professores XXXXXX e YYYYYY, sou grato pela orientação.

\item Ao secretário Halidaivson Stockhouse pela ajuda no andamento das coisas burocráticas do Curso.

\item Aos colegas ....

\item Aos demais colegas de graduação, pelas críticas e sugestões.

\item À minha família pelo apoio durante esta jornada.

\item À CAPES/CNPQ, pelo apoio financeiro.

\item À Deus.

\end{trivlist}


%
% O espaçamento entre linhas (ATENÇÃO A ESTA PARTE!)
%
%%%%%%%%%%%%%%%%%%%%%%%%%%%%%%%%%%%%%%%%%%%%%%%%%%%%%%%%%%%%%%%%%%%%%%%%%%%%
% PARA A VERSÃO FINAL:
% Deve ser usado espaçamento simples nas páginas de texto
\singlespacing
% PARA A QUALIFICAÇÃO E PARA A VERSÃO INICIAL:
% Deve ser usado espaçamento 1 1/2 nas páginas de texto
%\doublespacing
%%%%%%%%%%%%%%%%%%%%%%%%%%%%%%%%%%%%%%%%%%%%%%%%%%%%%%%%%%%%%%%%%%%%%%%%%%%%

% Resumo/Abstract
%
% ********** Resumo
%

% Usa-se \chapter*, e não \chapter, porque este "capítulo" não deve
% ser numerado.
% Na maioria das vezes, ao invés dos comandos LaTeX \chapter e \chapter*,
% deve-se usar as nossas versões definidas no arquivo comandos.tex,
% \mychapter e \mychapterast. Isto porque os comandos LaTeX têm um erro
% que faz com que eles sempre coloquem o número da página no rodapé na
% primeira página do capítulo, mesmo que o estilo que estejamos usando
% para numeração seja outro.
\mychapterast{Resumo}

Lorem ipsum dolor sit amet, consectetur adipiscing elit. Donec vehicula vitae lectus ut pretium. Vestibulum tristique leo eu purus vehicula ullamcorper. Nulla ut ultricies massa. Suspendisse eu neque pharetra, faucibus erat ac, pretium augue. Vivamus id euismod leo. Cras eget neque pellentesque, fringilla dolor eu, pretium libero. Mauris sed justo feugiat, varius ligula sed, posuere metus. Fusce lacus mi, molestie a rutrum id, scelerisque ut lacus. In hac habitasse platea dictumst. In vitae elit faucibus, molestie orci efficitur, consectetur neque. Ut placerat, augue eu pellentesque euismod, dui enim euismod elit, quis sollicitudin lectus lorem gravida mi. Donec ut leo pretium, finibus arcu in, tincidunt sem. Phasellus diam ante, pulvinar vel neque non, sagittis aliquam nibh. Praesent id condimentum nunc, quis interdum metus. Curabitur eget diam vitae enim consequat mollis quis dictum turpis.

\vspace{1.5ex}

{\bf Palavras-chave}: Processamento de texto, \LaTeX,
Preparação de Teses, Relatórios Técnicos.

%
% ********** Abstract
%
\mychapterast{Abstract}

Lorem ipsum dolor sit amet, consectetur adipiscing elit. Donec vehicula vitae lectus ut pretium. Vestibulum tristique leo eu purus vehicula ullamcorper. Nulla ut ultricies massa. Suspendisse eu neque pharetra, faucibus erat ac, pretium augue. Vivamus id euismod leo. Cras eget neque pellentesque, fringilla dolor eu, pretium libero. Mauris sed justo feugiat, varius ligula sed, posuere metus. Fusce lacus mi, molestie a rutrum id, scelerisque ut lacus. In hac habitasse platea dictumst. In vitae elit faucibus, molestie orci efficitur, consectetur neque. Ut placerat, augue eu pellentesque euismod, dui enim euismod elit, quis sollicitudin lectus lorem gravida mi. Donec ut leo pretium, finibus arcu in, tincidunt sem. Phasellus diam ante, pulvinar vel neque non, sagittis aliquam nibh. Praesent id condimentum nunc, quis interdum metus. Curabitur eget diam vitae enim consequat mollis quis dictum turpis.

\vspace{1.5ex}

{\bf Keywords}: Document Processing, \LaTeX, Thesis Preparation,
Technical Reports.


% Paginas introdutórias (com numeração romana)
\frontmatter

% Lista de conteúdo (sumário, gerado automaticamente)
% Aqui, e em todos os itens antes da introdução, o comando \phantomsection é utilizado.
% O seu uso é neecssário para auxiliar o pacote "hyperref" a fazer a referência correta
% dos links do sumário, das listas (de tabelas, figuras, algoritmos) com as páginas
% respectivas.
% Caso seja tirado, o "hyperref" irá apontar o link do sumário para o abstract, o link
% do sumário para a lista de figuras, o link das listas de figuras para a lista de tabelas,
% e assim sucessivamente.
\phantomsection
\addcontentsline{toc}{chapter}{Sumário}
\tableofcontents

% Lista de figuras (gerada automaticamente)
\cleardoublepage
\phantomsection
\addcontentsline{toc}{chapter}{Lista de Figuras}
\listoffigures

% Lista de tabelas (gerada automaticamente)
\cleardoublepage
\phantomsection
\addcontentsline{toc}{chapter}{Lista de Tabelas}
\listoftables

% Páginas do texto principal (com cabeçalho)
\mainmatter
\pagestyle{headings}

% Para facilitar a organização, foi criado um diretório para cada
% capítulo do documento, pois assim os arquivos das figuras ficam
% classificados por capítulos

% Cap. 1 - Introdução
%%
\mychapter{Introdução}
\label{Cap:Introducao}

Lorem ipsum dolor sit amet, consectetur adipiscing elit. Morbi tristique, orci mollis tincidunt dignissim, purus lectus molestie odio, vitae pharetra nisi sapien et justo. Fusce consequat et elit condimentum tincidunt. In eu venenatis tortor, quis lobortis tellus. Mauris tincidunt gravida ante. Pellentesque ante elit, lacinia interdum bibendum sed, cursus non orci. Ut et diam nec justo efficitur tincidunt. Fusce convallis facilisis varius. Integer tempor hendrerit maximus.

Donec pulvinar, libero semper pretium fringilla, arcu ligula iaculis urna, id faucibus ante massa a augue. Ut iaculis fermentum convallis. Vestibulum ac porttitor urna, sit amet dignissim dui. Curabitur sollicitudin tincidunt risus quis vulputate. In in leo ut mi mattis tempor sed non nunc. Sed et velit nibh. Duis mollis lectus risus, ut vestibulum purus semper eu. Donec rhoncus ex sed scelerisque dictum. Praesent sit amet neque accumsan, euismod sapien eget, tincidunt magna. Etiam nec diam ac elit feugiat aliquet. Pellentesque et neque metus. Nulla molestie libero sed ullamcorper mollis. In a risus quis felis eleifend dapibus at mollis odio. Curabitur nisl nibh, placerat ac quam nec, vulputate blandit mauris.

Vestibulum neque lacus, fringilla a urna quis, egestas tincidunt orci. Phasellus rutrum elit at mauris feugiat, non egestas tortor dictum. Aliquam faucibus, velit eu aliquam fermentum, lacus sem pellentesque nibh, et tristique urna nibh eget purus. Morbi vitae felis posuere, rutrum turpis quis, rutrum lorem. Nam lacinia cursus neque sit amet fermentum. Mauris eu mauris diam. Quisque lacinia consequat quam, at convallis urna dignissim eget. Proin sit amet varius sem, vel facilisis purus. Sed non blandit sapien. Sed auctor venenatis nibh eu ornare. Quisque euismod urna ligula, quis vestibulum sapien lacinia vehicula. Aenean sit amet vehicula felis, a imperdiet tellus. Fusce eu sem urna. Integer placerat nibh in purus consectetur mattis. Aliquam erat neque, tincidunt ac porta sit amet, egestas non lacus.



% Cap. 2 - Teoria (Fundamentação Teórica)
%%
%% Capítulo 2: Regras gerais de estilo
%%

\mychapter{Fundamentação Teórica}
\label{Cap:Teoria}

Lorem ipsum dolor sit amet, consectetur adipiscing elit. Morbi tristique, orci mollis tincidunt dignissim, purus lectus molestie odio, vitae pharetra nisi sapien et justo. Fusce consequat et elit condimentum tincidunt. In eu venenatis tortor, quis lobortis tellus. Mauris tincidunt gravida ante. Pellentesque ante elit, lacinia interdum bibendum sed, cursus non orci. Ut et diam nec justo efficitur tincidunt. Fusce convallis facilisis varius. Integer tempor hendrerit maximus.

Donec pulvinar, libero semper pretium fringilla, arcu ligula iaculis urna, id faucibus ante massa a augue. Ut iaculis fermentum convallis. Vestibulum ac porttitor urna, sit amet dignissim dui. Curabitur sollicitudin tincidunt risus quis vulputate. In in leo ut mi mattis tempor sed non nunc. Sed et velit nibh. Duis mollis lectus risus, ut vestibulum purus semper eu. Donec rhoncus ex sed scelerisque dictum. Praesent sit amet neque accumsan, euismod sapien eget, tincidunt magna. Etiam nec diam ac elit feugiat aliquet. Pellentesque et neque metus. Nulla molestie libero sed ullamcorper mollis. In a risus quis felis eleifend dapibus at mollis odio. Curabitur nisl nibh, placerat ac quam nec, vulputate blandit mauris.

% Cap. 3 - Trabalhos Relacionados
%%
%% Capítulo 2: Expressões matemáticas
%%

\mychapter{Trabalhos relacionados}
\label{Cap:TrabalhosRelacionados}
Lorem ipsum dolor sit amet, consectetur adipiscing elit. Morbi tristique, orci mollis tincidunt dignissim, purus lectus molestie odio, vitae pharetra nisi sapien et justo. Fusce consequat et elit condimentum tincidunt. In eu venenatis tortor, quis lobortis tellus. Mauris tincidunt gravida ante. Pellentesque ante elit, lacinia interdum bibendum sed, cursus non orci. Ut et diam nec justo efficitur tincidunt. Fusce convallis facilisis varius. Integer tempor hendrerit maximus.

Donec pulvinar, libero semper pretium fringilla, arcu ligula iaculis urna, id faucibus ante massa a augue. Ut iaculis fermentum convallis. Vestibulum ac porttitor urna, sit amet dignissim dui. Curabitur sollicitudin tincidunt risus quis vulputate. In in leo ut mi mattis tempor sed non nunc. Sed et velit nibh. Duis mollis lectus risus, ut vestibulum purus semper eu. Donec rhoncus ex sed scelerisque dictum. Praesent sit amet neque accumsan, euismod sapien eget, tincidunt magna. Etiam nec diam ac elit feugiat aliquet. Pellentesque et neque metus. Nulla molestie libero sed ullamcorper mollis. In a risus quis felis eleifend dapibus at mollis odio. Curabitur nisl nibh, placerat ac quam nec, vulputate blandit mauris.

\cite{art-solimaes03}, \cite{col-pinto00}, ..

% Cap. 4 - Problema
%%
%% Capítulo 4: Problema
%%

\mychapter{Problema}
\label{Cap:Problema}
Lorem ipsum dolor sit amet, consectetur adipiscing elit. Morbi tristique, orci mollis tincidunt dignissim, purus lectus molestie odio, vitae pharetra nisi sapien et justo. Fusce consequat et elit condimentum tincidunt. In eu venenatis tortor, quis lobortis tellus. Mauris tincidunt gravida ante. Pellentesque ante elit, lacinia interdum bibendum sed, cursus non orci. Ut et diam nec justo efficitur tincidunt. Fusce convallis facilisis varius. Integer tempor hendrerit maximus.

Donec pulvinar, libero semper pretium fringilla, arcu ligula iaculis urna, id faucibus ante massa a augue. Ut iaculis fermentum convallis. Vestibulum ac porttitor urna, sit amet dignissim dui. Curabitur sollicitudin tincidunt risus quis vulputate. In in leo ut mi mattis tempor sed non nunc. Sed et velit nibh. Duis mollis lectus risus, ut vestibulum purus semper eu. Donec rhoncus ex sed scelerisque dictum. Praesent sit amet neque accumsan, euismod sapien eget, tincidunt magna. Etiam nec diam ac elit feugiat aliquet. Pellentesque et neque metus. Nulla molestie libero sed ullamcorper mollis. In a risus quis felis eleifend dapibus at mollis odio. Curabitur nisl nibh, placerat ac quam nec, vulputate blandit mauris.

Vestibulum neque lacus, fringilla a urna quis, egestas tincidunt orci. Phasellus rutrum elit at mauris feugiat, non egestas tortor dictum. Aliquam faucibus, velit eu aliquam fermentum, lacus sem pellentesque nibh, et tristique urna nibh eget purus. Morbi vitae felis posuere, rutrum turpis quis, rutrum lorem. Nam lacinia cursus neque sit amet fermentum. Mauris eu mauris diam. Quisque lacinia consequat quam, at convallis urna dignissim eget. Proin sit amet varius sem, vel facilisis purus. Sed non blandit sapien. Sed auctor venenatis nibh eu ornare. Quisque euismod urna ligula, quis vestibulum sapien lacinia vehicula. Aenean sit amet vehicula felis, a imperdiet tellus. Fusce eu sem urna. Integer placerat nibh in purus consectetur mattis. Aliquam erat neque, tincidunt ac porta sit amet, egestas non lacus.

\section{Aliquam erat neque}
\label{Sec:flutuantes}

Vestibulum neque lacus, fringilla a urna quis, egestas tincidunt orci. Phasellus rutrum elit at mauris feugiat, non egestas tortor dictum. Aliquam faucibus, velit eu aliquam fermentum, lacus sem pellentesque nibh, et tristique urna nibh eget purus. Morbi vitae felis posuere, rutrum turpis quis, rutrum lorem. Nam lacinia cursus neque sit amet fermentum. Mauris eu mauris diam. Quisque lacinia consequat quam, at convallis urna dignissim eget. Proin sit amet varius sem, vel facilisis purus. Sed non blandit sapien. Sed auctor venenatis nibh eu ornare. Quisque euismod urna ligula, quis vestibulum sapien lacinia vehicula. Aenean sit amet vehicula felis, a imperdiet tellus. Fusce eu sem urna. Integer placerat nibh in purus consectetur mattis. Aliquam erat neque, tincidunt ac porta sit amet, egestas non lacus.


\section{Quisque euismod urna}

\begin{table}[htbp]
\begin{tabularx}{\linewidth}{|p{3cm}|X|l|} \hline
COLUNA p & COLUNA X & COLUNA l \\ \hline
Largura fixa (não depende do conteúdo) &
Expandível &
Ajustável \\ \hline
Alinhada no topo &
Alinhada à esquerda &
Alinhada à esquerda \\ \hline
\end{tabularx}
\\[0.5cm]
\begin{tabularx}{\linewidth}{|b{3cm}|C|r|} \hline
COLUNA b & COLUNA C (ver \texttt{comandos.tex}) & COLUNA r \\ \hline
Largura fixa (não depende do conteúdo) &
Expandível &
Ajustável \\ \hline
Alinhada na base &
Centralizada &
Alinhada à direita \\ \hline
\end{tabularx}
\caption{Tabelas com colunas de diferentes larguras e alinhamentos}
\label{Tab:larguracolunas}
\end{table}



% Cap. 5 - Implementação
%%
%% Capítulo 4: Figuras, gráficos e tabelas
%%

\mychapter{Implementação}
\label{Cap:Implementacao}
Lorem ipsum dolor sit amet, consectetur adipiscing elit. Morbi tristique, orci mollis tincidunt dignissim, purus lectus molestie odio, vitae pharetra nisi sapien et justo. Fusce consequat et elit condimentum tincidunt. In eu venenatis tortor, quis lobortis tellus. Mauris tincidunt gravida ante. Pellentesque ante elit, lacinia interdum bibendum sed, cursus non orci. Ut et diam nec justo efficitur tincidunt. Fusce convallis facilisis varius. Integer tempor hendrerit maximus.

\begin{algorithm}
%% \SetLine
\Entrada{$x$: vetores de valores; $y$ = $L(x)$; $p$: valor de entrada a ser calculado }
\Saida{$s$ = $L(p)$}
$n \leftarrow \mathtt{comprimento}(x)$\;
$s \leftarrow 0$\;
\Para {$i=1$ \Ate $n$} {
	$L \leftarrow 1$\;
	\Para {$j=1:1:n$} {
		\Se{$i \neq j$} {
			$L \leftarrow L* \left( \dfrac{p-x[j]}{x[i]-x[j]} \right) $
		}
	}
	$s \leftarrow s + L*y[i]$\;
}
\Retorna $s$\;
\caption{Algoritmo para interpolação de Lagrange.}
\label{algo:1}
\end{algorithm}

\begin{algorithm}
%% \SetLine
\Entrada{$a$: valor inicial; $b$: valor final; $n$: número de subintervalos (deve ser múltiplo de 2)  }
\tcc{A função a ser integrada é definida em uma função denominada \texttt{f}, fora do escopo deste algoritmo.}
\Saida{$I$ = integral de \texttt{f} entre $a$ e $b$}
$h \leftarrow$ $\dfrac{b-a}{n}$\;
$x[1] \leftarrow a$\;
$y[1] \leftarrow f(a)$\;
$I \leftarrow 0$\;
$k \leftarrow 2$\;
\Enqto {$k <= n$} {
	$x[i] \leftarrow x[i-1] + h$\;
	$y[i] \leftarrow f(x[i])$\;
	\eSe{$i \% 2 = 0$} {
		$I \leftarrow I + 4*y[i]$\;
	}
	{
		$I \leftarrow I + 2*y[i]$\;
	}
	$k = k+1$\;
}
$x[n+1] \leftarrow b$\;
$y[n+1] \leftarrow f(x[i+1])$\;
$I \leftarrow I + \dfrac{h}{3}*(I + y[n+1])$\;
\Retorna $I$\;
\caption{Algoritmo para a integração pelo primeiro método de Simpson.}
\label{algo:2}
\end{algorithm}



\section{Nulla molestie libero sed}
\label{Sec:expressoesMatematicas}

\begin{eqnarray} \label{eq:PDF:RSR}
  p \left( \gamma \right) & = & \frac{1}{2} \sqrt{\frac{M}{\gamma \bar{\gamma}_{b}}} \frac{1}{ \prod_{i=1}^M {\sqrt{\tilde{\gamma}_i}}}
  \int_0^{\sqrt{M \delta}} \int_0^{\sqrt{M \delta} - r_M } \cdots
  \int_0^{\sqrt{M \delta} - \sum_{i = 3}^M {r_i } } \nonumber \\
  & & p \left( {\frac{\sqrt{M \delta} - \sum_{i = 2}^M {r_i }}{\sqrt{\tilde{\gamma}_1}} ,
  \frac{r_2}{\sqrt{\tilde{\gamma}_2}} , \ldots ,\frac{r_M}{\sqrt{\tilde{\gamma}_M}} } \right)
  \, dr_2 \cdots dr_{M-1} \, dr_M
\end{eqnarray}
% sem linha em branco
ou:
% sem linha em branco
\begin{equation} \label{eq:TrCGI}
  T(r) = \frac{1}{f_m}
  \left( \frac{\pi}{2} \sum_{i=1}^M
  {\tilde{r}_i^2 \dot{\varsigma}_i^2}\right)^{-1/2}
  \frac
  {\begin{array}{ll}
  \int_0^{\rho \sqrt{M}} \int_0^{\rho \sqrt{M} - r_M } \cdots
  \int_0^{\rho \sqrt{M} - \sum_{i = 3}^M {r_i } } \int_0^{\rho \sqrt{M} -
  \sum_{i = 2}^M {r_i } }  \\
  p \left( {\frac{r_1}{\tilde{r}_1} ,
  \frac{r_2}{\tilde{r}_2} , \ldots ,\frac{r_M}{\tilde{r}_M} } \right)
  \, dr_1 \, dr_2 \cdots dr_{M-1} \, dr_M \\ \end{array}}
  {\begin{array}{ll}
  \int_0^{\rho \sqrt{M}} \int_0^{\rho \sqrt{M} - r_M } \cdots
  \int_0^{\rho \sqrt{M} - \sum_{i = 3}^M {r_i } } \\
  p \left( {\frac{\rho \sqrt{M} - \sum_{i = 2}^M {r_i }}{\tilde{r}_1} ,
  \frac{r_2}{\tilde{r}_2} , \ldots ,\frac{r_M}{\tilde{r}_M} } \right)
  \, dr_2 \cdots dr_{M-1} \, dr_M \\ \end{array}}
\end{equation}


\begin{equation}
y = g(x) = g(x_{PO}) + \left.\frac{dg}{dx}\right|_{x=x_{PO}}
\frac{(x-x_{PO})}{1!} + \left.\frac{d^2g}{dx^2}\right|_{x=x_{PO}}
\frac{(x-x_{PO})^2}{2!} + \cdots
\label{Eq:Taylor}
\end{equation}



% Cap. 6 - Experimentos e Resultados
%%
%% Capítulo 4: Figuras, gráficos e tabelas
%%

\mychapter{Experimentos e Resultados}
\label{Cap:ExperimentosResultados}

Lorem ipsum dolor sit amet, consectetur adipiscing elit. Morbi tristique, orci mollis tincidunt dignissim, purus lectus molestie odio, vitae pharetra nisi sapien et justo. Fusce consequat et elit condimentum tincidunt. In eu venenatis tortor, quis lobortis tellus. Mauris tincidunt gravida ante. Pellentesque ante elit, lacinia interdum bibendum sed, cursus non orci. Ut et diam nec justo efficitur tincidunt. Fusce convallis facilisis varius. Integer tempor hendrerit maximus.

Donec pulvinar, libero semper pretium fringilla, arcu ligula iaculis urna, id faucibus ante massa a augue. Ut iaculis fermentum convallis. Vestibulum ac porttitor urna, sit amet dignissim dui. Curabitur sollicitudin tincidunt risus quis vulputate. In in leo ut mi mattis tempor sed non nunc. Sed et velit nibh. Duis mollis lectus risus, ut vestibulum purus semper eu. Donec rhoncus ex sed scelerisque dictum. Praesent sit amet neque accumsan, euismod sapien eget, tincidunt magna. Etiam nec diam ac elit feugiat aliquet. Pellentesque et neque metus. Nulla molestie libero sed ullamcorper mollis. In a risus quis felis eleifend dapibus at mollis odio. Curabitur nisl nibh, placerat ac quam nec, vulputate blandit mauris.


% Cap. 7 - Conclusão
%%
%% Capítulo 5: Conclusões
%%

\mychapter{Conclusão}
\label{Cap:Conclusao}

Lorem ipsum dolor sit amet, consectetur adipiscing elit. Morbi tristique, orci mollis tincidunt dignissim, purus lectus molestie odio, vitae pharetra nisi sapien et justo. Fusce consequat et elit condimentum tincidunt. In eu venenatis tortor, quis lobortis tellus. Mauris tincidunt gravida ante. Pellentesque ante elit, lacinia interdum bibendum sed, cursus non orci. Ut et diam nec justo efficitur tincidunt. Fusce convallis facilisis varius. Integer tempor hendrerit maximus.

Donec pulvinar, libero semper pretium fringilla, arcu ligula iaculis urna, id faucibus ante massa a augue. Ut iaculis fermentum convallis. Vestibulum ac porttitor urna, sit amet dignissim dui. Curabitur sollicitudin tincidunt risus quis vulputate. In in leo ut mi mattis tempor sed non nunc. Sed et velit nibh. Duis mollis lectus risus, ut vestibulum purus semper eu. Donec rhoncus ex sed scelerisque dictum. Praesent sit amet neque accumsan, euismod sapien eget, tincidunt magna. Etiam nec diam ac elit feugiat aliquet. Pellentesque et neque metus. Nulla molestie libero sed ullamcorper mollis. In a risus quis felis eleifend dapibus at mollis odio. Curabitur nisl nibh, placerat ac quam nec, vulputate blandit mauris.

Vestibulum neque lacus, fringilla a urna quis, egestas tincidunt orci. Phasellus rutrum elit at mauris feugiat, non egestas tortor dictum. Aliquam faucibus, velit eu aliquam fermentum, lacus sem pellentesque nibh, et tristique urna nibh eget purus. Morbi vitae felis posuere, rutrum turpis quis, rutrum lorem. Nam lacinia cursus neque sit amet fermentum. Mauris eu mauris diam. Quisque lacinia consequat quam, at convallis urna dignissim eget. Proin sit amet varius sem, vel facilisis purus. Sed non blandit sapien. Sed auctor venenatis nibh eu ornare. Quisque euismod urna ligula, quis vestibulum sapien lacinia vehicula. Aenean sit amet vehicula felis, a imperdiet tellus. Fusce eu sem urna. Integer placerat nibh in purus consectetur mattis. Aliquam erat neque, tincidunt ac porta sit amet, egestas non lacus.

% Referências bibliogáficas (geradas automaticamente)
% Aqui, o comando \phantomsection é utilizado para auxiliar o pacote "hyperref" a fazer a
% referência correta dos links das referências bibliográficas com a página respectiva.
% Caso seja tirado, o "hyperref" irá apontar o link das referências bibliográficas para a
% última subseção da conclusão.
\phantomsection
\addcontentsline{toc}{chapter}{Referências bibliográficas}
\bibliography{bibliografia/bibliografia}

\appendix

%Apêndice A
%%
%% Apêndice 1: Requisitos funcionais detalhados
%%

\mychapter{Requisitos funcionais detalhados}
\label{apendice:requisitos_funcionais_detalhados}

\rowcolors{1}{white}{green!10}
\begin{longtable}{
  | >{\centering}m{0.10\textwidth} 
  | >{\centering}m{0.30\textwidth} 
  | m{0.60\textwidth} |}
  
    \hline
    \rowcolor{green!100}
    \textbf{Código} & \textbf{Requisito} & \textbf{Descrição} \\ \hline \hline
    \textbf{RF01}  &  \textbf{Cadastrar entidades} & \textbf{Permitir que o usuário cadastre novas entidades}  \\ \hline \hline
    RF01-01 & Cadastrar Vacinações              & Fornecer para preenchimento os campos obrigatórios aplicante, estabelecimento de aplicação, vacina, lote, informações do paciente (pré-cadastrado ou não), dose e data da aplicação da vacina e os campos opcionais campanha da vacinação que está sendo efetuada e data da próxima vacina \\ \hline
    RF01-02 & Cadastrar Campanhas               & Fornecer para preenchimento os campos obrigatórios título e datas de início de fim e o campo opcional descrição da campanha  \\ \hline
    RF01-03 & Cadastrar Estabelecimentos        & Fornecer para preenchimento os campos obrigatórios nome, CNES e endereço do estabelecimento de saúde \\ \hline
    RF01-04 & Cadastrar Vacinas                 & Fornecer para preenchimento os campos obrigatórios nome, fabricante e código da vacina \\ \hline
    RF01-05 & Cadastrar Lotes de Vacina         & Fornecer para preenchimento os campos obrigatórios código, nome e quantidade de doses da vacina no lote \\ \hline
    RF01-06 & Cadastrar Pacientes               & Fornecer para preenchimento os campos obrigatórios CNS, CPF, nome, data de nascimento, localidade, categoria prioritária e condição maternal e os campos opcionais sexo e nomes do pai e da mãe do paciente \\ \hline
    RF01-07 & Cadastrar Aplicantes              & Fornecer para preenchimento os campos obrigatórios CNS, CPF, nome, localidade e estabelecimento de saúde de atuação e os campos opcionais data de nascimento, sexo e nomes do pai e da mãe do paciente \\ \hline
    RF01-08 & Cadastrar Localidades             & Fornecer para preenchimento os campos obrigatórios código do IBGE, nome, cidade e estado da localidade \\ \hline
    RF01-09 & Cadastrar Grupos Prioritários     & Fornecer para preenchimento o campo obrigatório código e os campos opcionais nome e descrição do grupo prioritário \\ \hline
    RF01-10 & Cadastrar Categorias Prioritárias & Fornecer para preenchimento os campos obrigatórios* código e grupo prioritário pertencente e os campos opcionais nome e descrição da categoria prioritária \\ \hline
    RF01-11 & Botão de salvamento               & Fornecer para preenchimento os campos obrigatórios* botão para salvamento dos dados cadastrados em todos os formulários \\ \hline

  \hiderowcolors
  \caption{Requisitos funcionais da aplicação Nurse: detalhes do requisito RF01}
  \label{tab:rf01_detalhe}
\end{longtable}

% \begin{table}[ht!]
%   \centering
%   {\rowcolors{0}{white}{green!10}
%   \begin{tabularx}{\textwidth}{
%     | >{\centering\arraybackslash}m{0.10\textwidth} 
%     | >{\centering\arraybackslash}X 
%     | >{\raggedright\arraybackslash}X | }
%     \hline
%     \rowcolor{green!100}
%     \textbf{Código} & \textbf{Requisito} & \textbf{Descrição} \\ \hline \hline
%     \textbf{RF01}  &  \textbf{Cadastrar entidades} & \textbf{Permitir que o usuário cadastre novas entidades}  \\ \hline \hline
%     RF01-01 & Cadastrar Vacinações              & Fornecer para preenchimento os campos obrigatórios aplicante, estabelecimento de aplicação, vacina, lote, informações do paciente (pré-cadastrado ou não), dose e data da aplicação da vacina e os campos opcionais campanha da vacinação que está sendo efetuada e data da próxima vacina \\ \hline
%     RF01-02 & Cadastrar Campanhas               & Fornecer para preenchimento os campos obrigatórios título e datas de início de fim e o campo opcional descrição da campanha  \\ \hline
%     RF01-03 & Cadastrar Estabelecimentos        & Fornecer para preenchimento os campos obrigatórios nome, CNES e endereço do estabelecimento de saúde \\ \hline
%     RF01-04 & Cadastrar Vacinas                 & Fornecer para preenchimento os campos obrigatórios nome, fabricante e código da vacina \\ \hline
%     RF01-05 & Cadastrar Lotes de Vacina         & Fornecer para preenchimento os campos obrigatórios código, nome e quantidade de doses da vacina no lote \\ \hline
%     RF01-06 & Cadastrar Pacientes               & Fornecer para preenchimento os campos obrigatórios CNS, CPF, nome, data de nascimento, localidade, categoria prioritária e condição maternal e os campos opcionais sexo e nomes do pai e da mãe do paciente \\ \hline
%     RF01-07 & Cadastrar Aplicantes              & Fornecer para preenchimento os campos obrigatórios CNS, CPF, nome, localidade e estabelecimento de saúde de atuação e os campos opcionais data de nascimento, sexo e nomes do pai e da mãe do paciente \\ \hline
%     RF01-08 & Cadastrar Localidades             & Fornecer para preenchimento os campos obrigatórios código do IBGE, nome, cidade e estado da localidade \\ \hline
%     RF01-09 & Cadastrar Grupos Prioritários     & Fornecer para preenchimento o campo obrigatório código e os campos opcionais nome e descrição do grupo prioritário \\ \hline
%     RF01-10 & Cadastrar Categorias Prioritárias & Fornecer para preenchimento os campos obrigatórios* código e grupo prioritário pertencente e os campos opcionais nome e descrição da categoria prioritária \\ \hline
%     RF01-11 & Botão de salvamento               & Fornecer para preenchimento os campos obrigatórios* botão para salvamento dos dados cadastrados em todos os formulários \\ \hline
%   \end{tabularx}}
% \caption{Requisitos funcionais da aplicação Nurse: detalhes do requisito RF01}
% \label{tab:rf01_detalhe}
% \end{table}

\begin{table}[ht!]
  \centering
  {\rowcolors{0}{white}{green!10}
  \begin{tabularx}{\textwidth}{
    | >{\centering\arraybackslash}m{0.10\textwidth} 
    | >{\centering\arraybackslash}m{0.30\textwidth} 
    | >{\raggedright\arraybackslash}X | }
    \hline
    \rowcolor{green!100}
    \textbf{Código} & \textbf{Requisito} & \textbf{Descrição} \\ \hline \hline
    \textbf{RF02}  &  \textbf{Visualizar entidades cadastradas} & \textbf{Permitir que o usuário visualize as entidades cadastradas e seus detalhes}  \\ \hline \hline
    RF02-01 & Visualizar Vacinações               & Apresentar CNS, nome e grupo do paciente e vacina aplicada \\ \hline
    RF02-02 & Visualizar Campanhas                & Apresentar título, datas de início de fim e status da campanha \\ \hline
    RF02-03 & Visualizar Estabelecimentos         & Apresentar nome, CNES e endereço do estabelecimento de saúde \\ \hline
    RF02-04 & Visualizar Vacinas                  & Apresentar nome, fabricante e código da vacina \\ \hline
    RF02-05 & Visualizar Lotes de Vacina          & Apresentar código, nome e quantidade de doses da vacina no lote \\ \hline
    RF02-06 & Visualizar Pacientes                & Apresentar CNS, nome, categoria prioritária e condição maternal do paciente \\ \hline
    RF02-07 & Visualizar Aplicantes               & Apresentar CNS, nome e estabelecimento de saúde do profissional \\ \hline
    RF02-08 & Visualizar Localidades              & Apresentar código do IBGE, nome e endereço da localidade \\ \hline
    RF02-09 & Visualizar Grupos Prioritários      & Apresentar código, nome e descrição do grupo prioritário \\ \hline
    RF02-10 & Visualizar Categorias Prioritárias  & Apresentar código, nome e descrição da categoria prioritária \\ \hline
  \end{tabularx}}
\caption{Requisitos funcionais da aplicação Nurse: detalhes do requisito RF02}
\label{tab:rf02_detalhe}
\end{table}

\begin{table}[ht!]
  \centering
  {\rowcolors{0}{white}{green!10}
  \begin{tabularx}{\textwidth}{
    | >{\centering\arraybackslash}m{0.10\textwidth} 
    | >{\centering\arraybackslash}m{0.30\textwidth} 
    | >{\raggedright\arraybackslash}X | }
    \hline
    \rowcolor{green!100}
    \textbf{Código} & \textbf{Requisito} & \textbf{Descrição} \\ \hline \hline
    \textbf{RF04}  &  \textbf{Gerar tabela de vacinações} & \textbf{Permitir que o usuário gere uma
    tabela de vacinações para
    exportação}  \\ \hline \hline
    RF04-01 & Escolher período de tempo               & Permitir que o usuário escolha o período desejado para coleta dos dados sobre vacinação a serem exportados \\ \hline
    RF04-02 & Escolher o que fazer com o arquivo gerado                     & Permitir que o usuário escolha entre abrir o arquivo com a planilha gerada ou exportá-la utilizando para isso alguma aplicação externa compatível com a escolha, a depender do caso \\ \hline
  \end{tabularx}}
\caption{Requisitos funcionais da aplicação Nurse: detalhes do requisito RF04}
\label{tab:rf04_detalhe}
\end{table}

\end{document}
